\documentclass[12pt,a4paper]{article}

% Paquetes esenciales para tipografía y formato académico
\usepackage[utf8]{inputenc}
\usepackage[T1]{fontenc}
\usepackage{lmodern} % Tipografía Latin Modern - alta legibilidad
\usepackage{microtype} % Mejoras tipográficas finas
\usepackage{geometry} % Control de márgenes
\usepackage{setspace} % Control de interlineado
\usepackage{graphicx} % Para imágenes
\usepackage{float} % Control de posición de figuras
\usepackage{caption} % Mejoras en captions
\usepackage{subcaption} % Subfiguras
\usepackage{booktabs} % Tablas profesionales
\usepackage{array} % Mejoras en tablas
\usepackage{multirow} % Celdas múltiples en tablas
\usepackage{tabularx} % Tablas con ancho ajustable
\usepackage{longtable} % Tablas multipágina
\usepackage{adjustbox} % Ajuste de cajas y tablas
\usepackage{xcolor} % Colores
\usepackage{hyperref} % Enlaces y PDF interactivo
\usepackage{tikz} % Para diagramas
\usepackage{fancyhdr} % Encabezados y pies de página
\usepackage{enumitem} % Listas mejoradas
\usepackage{titlesec} % Títulos personalizados
\usetikzlibrary{matrix, positioning, arrows.meta}

% Configuración de márgenes optimizada para formato académico
\geometry{
    a4paper,
    left=2.5cm,      % Margen izquierdo estándar académico
    right=2.5cm,     % Margen derecho
    top=2.5cm,       % Margen superior
    bottom=2.5cm,    % Margen inferior
    headheight=15pt,
    headsep=20pt,
    footskip=30pt
}

% Configuración de interlineado para mejor legibilidad
\renewcommand{\baselinestretch}{1.3} % 1.3 de interlineado (estándar académico)

% Configuración de colores para alto contraste
\definecolor{darkblue}{RGB}{0,51,102}
\definecolor{lightgray}{RGB}{245,245,245}
\definecolor{fortaleza}{RGB}{76,175,80}
\definecolor{oportunidad}{RGB}{33,150,243}
\definecolor{debilidad}{RGB}{244,67,54}
\definecolor{amenaza}{RGB}{255,152,0}

% Configuración de hyperref para accesibilidad
\hypersetup{
    colorlinks=true,
    linkcolor=darkblue,
    filecolor=darkblue,
    urlcolor=darkblue,
    citecolor=darkblue,
    bookmarks=true,
    bookmarksopen=true,
    bookmarksnumbered=true,
    pdfstartview=FitH
}

% Configuración de fuentes y tamaños
\setlength{\parindent}{0pt} % Sin sangría
\setlength{\parskip}{8pt} % Espacio entre párrafos

% Configuración de títulos estilo académico
\titleformat{\section}{\Large\bfseries\scshape}{\thesection}{1em}{}
\titleformat{\subsection}{\large\bfseries}{\thesubsection}{1em}{}
\titleformat{\subsubsection}{\normalsize\bfseries}{\thesubsubsection}{1em}{}

% Espaciado entre secciones
\titlespacing*{\section}{0pt}{20pt}{10pt}
\titlespacing*{\subsection}{0pt}{15pt}{8pt}
\titlespacing*{\subsubsection}{0pt}{12pt}{6pt}

% Configuración de encabezado y pie de página
\pagestyle{fancy}
\fancyhf{}
\fancyhead[L]{\small\textsc{EPRO Grupo 1}}
\fancyhead[R]{\small\textsc{Universidad Escuela Colombiana de Ingeniería}}
\fancyfoot[C]{\small\thepage}
\renewcommand{\headrulewidth}{0.4pt}
\renewcommand{\footrulewidth}{0pt}

% Comandos personalizados para consistencia
\newcommand{\sectiontitle}[1]{\section{#1}}
\newcommand{\subsectiontitle}[1]{\subsection{#1}}
\newcommand{\subsubsectiontitle}[1]{\subsubsection{#1}}
\newcommand{\tablecaption}[1]{\caption{#1}}

% Configuración para tablas
\renewcommand{\arraystretch}{1.2} % Espacio en celdas de tabla
\setlength{\tabcolsep}{8pt} % Espacio entre columnas

% Configuración para figuras
\captionsetup[figure]{font=small,labelfont=bf}
\captionsetup[table]{font=small,labelfont=bf}

% Configuración de listas
\setlist[itemize]{leftmargin=*,itemsep=5pt,topsep=5pt}
\setlist[enumerate]{leftmargin=*,itemsep=5pt,topsep=5pt}

% Comando para tablas anchas
\newcommand{\adjusttable}[1]{%
  \begin{adjustbox}{width=\textwidth,center}
    #1
  \end{adjustbox}
}

% Información del documento
\title{\textbf{MATRIZ DOFA DEL PROYECTO:}\\
MONTAJE DE UNA EMPRESA DE RESERVAS CON PLANIFICACIÓN DE VIAJE, A\\
TRAVÉS DE UNA APLICACIÓN EN COLOMBIA}
\author{EPRO Grupo 1\\
    \small Santiago Botero García\\
    \small Juana Lozano Chaves\\
    \small Juan Pablo Salamanca Mojica\\
    \small Sofia Velandia Cifuentes\\
    \small Laura Alejandra Venegas Piraban\\
    \vspace{0.5cm}
    \small Profesor: Daniel Salazar Ferro}
\date{BOGOTÁ, COLOMBIA}

\begin{document}

% Página de título
\begin{titlepage}
    \centering
    \vspace*{2cm}
    
    {\Large\textbf{UNIVERSIDAD ESCUELA COLOMBIANA DE INGENIERÍA JULIO GARAVITO}}\\[1cm]
    
    {\Large\textbf{MATRIZ DOFA DEL PROYECTO:}}\\[0.5cm]
    {\Large MONTAJE DE UNA EMPRESA DE RESERVAS CON PLANIFICACIÓN DE VIAJE, A}\\[0.3cm]
    {\Large TRAVÉS DE UNA APLICACIÓN EN COLOMBIA}\\[2cm]
    
    {\large\textbf{EPRO GRUPO 1}}\\[1.5cm]
    
    \textbf{Estudiantes:}\\[0.5cm]
    JUAN PABLO SALAMANCA MOJICA\\
    JUANA LOZANO CHAVES\\
    LAURA ALEJANDRA VENEGAS\\
    SANTIAGO BOTERO GARCÍA\\
    SOFÍA VELANDIA CIFUENTES\\[1.5cm]
    
    \textbf{Profesor:}\\[0.5cm]
    Daniel Salazar Ferro\\[2cm]
    
    \textbf{BOGOTÁ, COLOMBIA}
    
    \vfill
\end{titlepage}

% Tabla de contenidos
\tableofcontents
\newpage

% Contenido principal
\sectiontitle{MATRIZ DOFA}

\begin{table}[H]
\centering
\tiny
\begin{tabular}{|p{2.8cm}|p{4.2cm}|p{4.2cm}|}
\hline
\textbf{Factores Internos/Externos} & \textbf{OPORTUNIDADES (O)} & \textbf{AMENAZAS (A)} \\
\hline
\textbf{FORTALEZAS (F)} & 
\textbf{Estrategias FO} & 
\textbf{Estrategias FA} \\
1. App única (reservas + planificación + red social) & 
\textbf{P1 (F1O1,2,8):} Montaje de una empresa de reservas con planificación de viaje y conexión entre viajeros mediante app & 
\textbf{P1 (A7,9):} Montaje de la app como alternativa a OTA's y plataformas informales \\
2. Seguridad y protección de datos & 
\textbf{P2 (F1O6):} Implementación de un portafolio de destinos top en Colombia & 
\textbf{P2 (F3A2,7):} Fortalecer cumplimiento normativo y formalidad para generar confianza \\
3. Enfoque en destinos emergentes & 
\textbf{P3 (F2O3):} Implementación de sistema de pagos digitales optimizado & 
\textbf{P3 (F2A7,9):} Ofrecer acompañamiento local como valor diferencial frente a plataformas globales \\
& 
\textbf{P1 adicional:} Uso de incentivos del MinCIT para el desarrollo del producto & 
\\
\hline
\textbf{DEBILIDADES (D)} & 
\textbf{Estrategias DO} & 
\textbf{Estrategias DA} \\
1. Recursos limitados & 
\textbf{P1:} Uso de incentivos gubernamentales para compensar recursos limitados & 
\textbf{P1 (A2,7,9D2,4,5):} Estrategias para posicionamiento de marca y generación de confianza \\
2. Marca sin reconocimiento & 
\textbf{P2 (D2O4):} Campaña publicitaria en aeropuertos con destinos top & 
\textbf{P2 (A1,5,6D1,5):} Implementación de trabajo remoto, software en la nube y contratación por servicios para reducir costos \\
3. Alto mantenimiento & 
\textbf{P3:} Aprovechar la digitalización del turismo para posicionar la marca & 
\textbf{P3 (A3,4,10D3,5):} Desarrollo de app funcional sin conexión (mapas, tiquetes, datos) \\
4. Poca experiencia & 
 & 
\\
\hline
\end{tabular}
\caption{Matriz DOFA del proyecto de turismo digital}
\end{table}

\sectiontitle{ANÁLISIS DETALLADO DE FACTORES}

\subsectiontitle{Fortalezas (Internas)}

\begin{enumerate}
    \item \textbf{Única aplicación en el mercado local que:}
    \begin{itemize}
        \item Integra reservas, planificación algorítmica y red social de viajeros.
        \item Reduce dependencia de intermediarios y fortalece el turismo nacional.
        \item Integra de múltiples actores (hoteles, guías, transporte).
    \end{itemize}
    \item \textbf{Protección de datos y seguridad ya integrados.}
    \item \textbf{Desarrollo en destinos emergentes.}
\end{enumerate}

\subsectiontitle{Debilidades}

\begin{enumerate}
    \item \textbf{Baja disponibilidad de recursos}
    \item \textbf{Marca nueva sin reconocimiento.}
    \item \textbf{Necesidad de mantenimiento constante.}
    \item \textbf{Poca experiencia en el sector empresarial y turismo.}
\end{enumerate}

\subsectiontitle{Oportunidades}

\textbf{Entorno PESTA (Cifras a 2026)}

\begin{enumerate}
    \item \textbf{(P) Político:} El Gobierno Nacional proyecta un cierre de 2025 con 21,7 millones de viajeros (MinCIT) (Eje de transición económica)
    \item \textbf{(E) Económico:} Crecimiento de un 6,4\% en reservas aéreas internacionales a Colombia en 2026.
    \item \textbf{(S) Social:} El comercio electrónico en Colombia registró 684 millones de transacciones en 2025 (crecimiento del 19,9\%).
    \item \textbf{(T) Tecnológico:} Mayor conectividad aérea en aeropuertos secundarios: Santa Marta (+11,9\%) y Medellín (+5,6\%), facilitando el acceso a destinos emergentes.
    \item \textbf{(A) Ambiental:} La demanda de destinos con baja huella de carbono aumentó al 11,5\% en zonas rurales.
\end{enumerate}

\textbf{Sector PORTER}

\begin{enumerate}
    \setcounter{enumi}{5}
    \item \textbf{(RC) Rivalidad Competitiva:} Chocó, La Guajira y Vichada son los destinos con mayor crecimiento de visitantes extranjeros (+4,3\%).
    \item \textbf{(NP) Nuevos Participantes:} El 30,5\% de la población está en situación de vulnerabilidad, lo que permite integrar a pequeños prestadores locales a la cadena de valor digital.
    \item \textbf{Digitalización de la cadena turística tradicional.}
\end{enumerate}

\subsectiontitle{Amenazas}

\textbf{Entorno PESTA:}

\begin{enumerate}
    \item \textbf{(E) Económico:} El mayor gasto diario en Bogotá y Antioquia, más aranceles e impuestos, encarece el destino y limita a la clase media.
    \item \textbf{(S) Social:} Persiste desconfianza en pagos digitales; violencia y orden público afectan la intención de viaje.
    \item \textbf{(T) Tecnológico:} Limitada tecnología en el ``último kilómetro'' dificulta reservas, pagos y gestión en destinos de naturaleza.
    \item \textbf{(A) Ambiental:} Cambio climático y frentes de frío afectan vías al Pacífico y Amazonía, generando cierres y cancelaciones.
    \item \textbf{(E) Económico:} La volatilidad del dólar eleva costos operativos y reduce márgenes del sector turístico.
    \item \textbf{(P) Político:} Incertidumbre política y elecciones frenan inversión y consumo en turismo.
\end{enumerate}

\textbf{Sector PORTER:}

\begin{enumerate}
    \setcounter{enumi}{6}
    \item \textbf{(RC) Rivalidad Competitiva:} OTA's como Booking.com y Airbnb dominan el mercado y presionan comisiones.
    \item \textbf{(NC) Negociación con Clientes:} Usuarios comparan precios en tiempo real, aumentando sensibilidad y presión sobre tarifas.
    \item \textbf{(APS) Amenaza de Productos Sustitutos:} Instagram y TikTok facilitan reservas directas, elevando competencia informal.
    \item \textbf{(T) Tecnológico:} Infraestructura deficiente en vías y servicios limita competitividad y experiencia turística.
\end{enumerate}

\sectiontitle{ESTRATEGIAS DEFINIDAS}

\subsectiontitle{Estrategias de Desarrollo del Producto}

\begin{enumerate}
    \item \textbf{P1: Montaje de una empresa de reservas con planificación de viaje y conexión entre viajeros a través de una app en Colombia.}
    \begin{itemize}
        \item Aprovecha fortalezas únicas de integración (F1O1,2,8)
        \item Responde a amenazas de competencia global (A7;9)
        \item Mitiga debilidades de experiencia mediante tecnología (A2;7;9D2;4;5)
    \end{itemize}
    
    \item \textbf{P2: Fortalecer el cumplimiento normativo y formalidad para aumentar la confianza del consumidor (23,8\%) y diferenciarse de plataformas globales.}
    \begin{itemize}
        \item Capitaliza fortaleza de seguridad integrada (F3A2;7)
        \item Responde a desconfianza en pagos digitales (A2;7)
    \end{itemize}
    
    \item \textbf{P3: Ofrecer conexión humana y acompañamiento local como valor diferencial frente a OTA's y reservas informales en redes sociales.}
    \begin{itemize}
        \item Aprovecha fortaleza de seguridad (F2A7;9)
        \item Diferencia de competencia informal (A7;9)
    \end{itemize}
\end{enumerate}

\subsectiontitle{Estrategias de Desarrollo de Mercado}

\begin{enumerate}
    \setcounter{enumi}{3}
    \item \textbf{P2: Implementación de un portafolio de destinos top en Colombia.}
    \begin{itemize}
        \item Aprovecha fortaleza de integración múltiple (F1O6)
        \item Responde a crecimiento en destinos emergentes
    \end{itemize}
\end{enumerate}

\subsectiontitle{Estrategias de Penetración de Mercado}

\begin{enumerate}
    \setcounter{enumi}{4}
    \item \textbf{P2: Montaje de una campaña publicitaria a través de vallas electrónicas de destino top en aeropuertos de Colombia.}
    \begin{itemize}
        \item Mitiga debilidad de marca nueva (D2O4)
        \item Aprovecha mayor conectividad aérea (O4)
    \end{itemize}
    
    \item \textbf{P2: Implementación de una empresa de turismo digital con trabajo remoto, uso de software en la nube y contratación por servicio en Colombia.}
    \begin{itemize}
        \item Mitiga debilidad de recursos limitados (A1;5;6D1;5)
        \item Responde a volatilidad económica y política
    \end{itemize}
    
    \item \textbf{P3: Desarrollo de una aplicación de reservas que permita guardar tiquetes, mapas y datos del viaje sin conexión a internet en Colombia.}
    \begin{itemize}
        \item Mitiga debilidad de experiencia (A3;4;10D3;5)
        \item Responde a limitaciones tecnológicas en último kilómetro
    \end{itemize}
\end{enumerate}

\sectiontitle{HOJA DE RUTA}

Primero se implementará una \textbf{Estrategia de desarrollo del Producto}, con el proyecto: Montaje de una empresa de reservas con planificación de viaje y conexión entre viajeros a través de una app en Colombia. 

Posteriormente se implementará una \textbf{Estrategia de Desarrollo del Mercado}, con el proyecto: Implementación de un portafolio de destinos top en Colombia. 

Finalmente, se implementará una \textbf{Estrategia de Penetración del Mercado}, con el proyecto: Montaje de una campaña publicitaria a través de vallas electrónicas de destino top en aeropuertos de Colombia.

\sectiontitle{REFERENCIAS}

\begin{itemize}
    \item MinCIT (Ministerio de Comercio, Industria y Turismo) - Informes de comportamiento turístico y Plan Sectorial de Turismo.
    \item ANATO (Asociación Colombiana de Agencias de Viajes y Turismo) y ForwardKeys (analítica de datos aéreos).
    \item Cámara Colombiana de Comercio Electrónico (CCCE) - Informe Anual 2025.
    \item DANE (Departamento Administrativo Nacional de Estadística) - Encuesta Nacional de Calidad de Vida.
    \item Fedesarrollo - Encuesta de Opinión del Consumidor (Enero 2026).
    \item Aeronáutica Civil de Colombia (Aerocivil) - Estadísticas de transporte aéreo.
\end{itemize}

\end{document}
